% ============================================================
% MISE EN PAGE — VARIANTE « UNE COLONNE » (grands groupes / ATS)
% ============================================================
% Marges resserrées mais au-dessus du seuil de rognage des imprimantes
% bureautiques (~5 mm) : c'est le plancher tenable pour rester sur une page.
\geometry{left=1.3cm, right=1.3cm, top=0.7cm, bottom=0.6cm}

% Titres de section pleine largeur : un filet horizontal continu structure la
% lecture verticale, seul repère de hiérarchie une fois la sidebar supprimée.
\titlespacing*{\section}{0pt}{4pt}{2pt}

% Séparation entre deux entrées. Volontairement plus serrée que ce que la
% pleine largeur appellerait : la densité prime ici, et les filets de section
% suffisent à découper la lecture.
\setlength{\cvitemsep}{2pt}

% Interlignage resserré. ATTENTION : ce fichier est calibré au plus juste — le
% CV tient sur une page, mais sans marge. Détendre l'interlignage à 0,93, ou
% remonter une marge verticale d'un millimètre, suffit à le faire déborder
% (vérifié). Toute puce ajoutée doit donc en remplacer une autre, et chaque
% puce doit tenir sur UNE ligne (cf. l'en-tête de sections/experience.tex).
% Contrôle : ./build.sh onecolumn puis « pdfinfo cv-onecolumn.pdf | grep Pages ».
\linespread{0.92}\selectfont

% ------------------------------------------------------------
% FILIGRANE : ICÔNE OMARCHY CENTRÉE SUR CHAQUE PAGE
% ------------------------------------------------------------
% eso-pic dépose un contenu en fond de CHAQUE page au moment du « shipout » :
% le filigrane est donc répété automatiquement, sans avoir à savoir combien de
% pages fait le CV ni où tombent les coupures. Contrairement à un overlay TikZ
% « remember picture », cela ne demande pas de seconde passe de compilation.
\usepackage{eso-pic}

% \omarchywatermarksize : côté du filigrane (carré). Nettement plus grand que
% l'icône de la sidebar : au centre d'une page blanche pleine largeur, une
% petite icône passerait pour une tache plutôt que pour un fond.
\newlength{\omarchywatermarksize}
\setlength{\omarchywatermarksize}{11cm}

% Même parti pris que dans la sidebar : plutôt que « opacity », on mélange la
% couleur au blanc du papier. Le rendu est identique sur un fond uni, mais la
% couleur reste opaque et ordinaire — pas de canal de transparence, donc aucun
% risque à l'impression ni dans les lecteurs PDF qui aplatissent mal les
% calques. C'est aussi ce qui garantit que le texte par-dessus reste noir franc.
% \omarchywatermarkmix : part de maincolor, en pourcentage (0 = invisible).
% Volontairement très bas : le filigrane doit se deviner, pas se lire.
\newcommand{\omarchywatermarkmix}{4}

% \makebox(0,0){...} centre la boîte sur le point visé (le centre de la page),
% en largeur comme en hauteur : c'est l'idiome « picture » attendu par eso-pic.
% Fond (« BG ») et non premier plan : le texte du CV passe par-dessus.
\AddToShipoutPictureBG{%
  \AtPageCenter{%
    \makebox(0,0){%
      \omarchyicon{\omarchywatermarksize}{maincolor!\omarchywatermarkmix!white}%
    }%
  }%
}

% ------------------------------------------------------------
% COMMANDES SPÉCIFIQUES À LA VARIANTE UNE COLONNE
% ------------------------------------------------------------

% Puce textuelle plutôt qu'icône FontAwesome. Visuellement quasi identique,
% mais c'est un vrai caractère du flux de texte : pdftotext — et donc les ATS,
% cible de cette variante — le restituent à sa place, au lieu de rejeter les
% icônes en fin de bloc où elles viennent se coller au texte suivant.
\renewcommand{\cvbullet}{\textcolor{maincolor}{\small\textbullet}}

% Sous-puces de détail composées un cran plus petit qu'en variante sidebar :
% sur 18 cm de largeur, \footnotesize reste parfaitement lisible, et la plupart
% des détails passent alors sur une seule ligne au lieu de deux.
\renewcommand{\cvsub}[1]{{\footnotesize\color{detailcolor}\itshape #1}}

% Ligne de compétences : \cvskillline{Catégorie}{éléments}
% Le libellé est aligné sur une largeur fixe pour que toutes les catégories
% forment une colonne nette, le reste coulant en paragraphe pendu.
% Les \strut fixent la même hauteur de première ligne dans les deux boîtes :
% sans eux, les deux colonnes n'ont pas la même hauteur d'ascendante et
% « minipage[t] » cale les libellés une demi-ligne au-dessus des valeurs.
\newcommand{\cvskillline}[2]{%
  \par\noindent
  \begin{minipage}[t]{2.6cm}\raggedright\strut\textbf{\color{maincolor}#1}\end{minipage}%
  \hspace{4pt}%
  \begin{minipage}[t]{\dimexpr\linewidth-2.6cm-4pt\relax}\raggedright\strut\color{secondcolor}#2\end{minipage}%
  \par\vspace{3pt}%
}

% Séparateur des lignes de contact / d'énumérations en ligne.
\newcommand{\cvsep}{\hspace{6pt}\textcolor{maincolor}{\small\textbullet}\hspace{6pt}}
